\chapter{Misbehavior, Crises, Regulation and Self Regulation}

This lecture begins where the previous discussion of behavioral finance leaves off. If financial markets are exposed to human weakness, manipulation, overconfidence, and bad character, then we must ask what institutions do about these things. The answer is regulation, but regulation arrives in two different registers at once. It is partly moral and behavioral, aimed at fraud, abuse, and exploitation; it is also technical, aimed at keeping the financial system from breaking down. The chapter follows that progression closely: from \emph{too big to fail}, to prudential regulation, to the need for referees, and then upward through the increasingly large scales at which financial rules are made and enforced.

\section{Regulation After Behavioral Finance}

We begin with the lecturer's opening move: regulation is not an external add-on to finance, but a continuation of the study of human behavior in markets. If people exploit one another's weaknesses, then rules are needed to limit manipulation; if markets generate monopolies and externalities, then rules are also needed to make the system work at all.

The lecture chooses one example very quickly, because it forces us to think at the system level rather than only at the level of isolated wrongdoing. That example is \emph{too big to fail}. A large financial firm may fail, but its failure is not treated like the failure of an ordinary small firm. Governments repeatedly rescue such firms because they fear that letting them collapse would bring down other institutions with them.

This can be written schematically as
\begin{align}
P(\text{bailout}\mid \text{large-firm failure})
>
P(\text{bailout}\mid \text{small-firm failure}),
\end{align}
and once we write it in this way, the incentive problem is already visible. The large firm is effectively given a contingent guarantee.

The lecture's causal chain is then
\begin{align}
\text{implicit guarantee}
\;\Rightarrow\;
\text{greater risk-taking}
\;\Rightarrow\;
\text{greater systemic risk}.
\end{align}

This is not a formal theorem in the lecture, but it is the central mechanism. A bailout expectation does not merely soften one future failure; it changes present behavior.

\section{Too Big to Fail and Prudential Regulation}

The lecturer next introduces the distinction between \emph{microprudential} and \emph{macroprudential} regulation. ``Prudential'' here means regulation concerned with prudent standards of business conduct. But prudent with respect to what? The answer splits in two directions:
\begin{align}
\text{Microprudential regulation}
&\longrightarrow
\text{protection of firms, households, and small investors},\\
\text{Macroprudential regulation}
&\longrightarrow
\text{protection of the financial system as a system}.
\end{align}

The distinction matters because the same financial event can be described in both languages, but the lecture wants us to keep the objects of protection separate. Misleading mortgage terms, abusive borrowing arrangements, and small-investor deception are microprudential concerns. \emph{Too big to fail} is the paradigm macroprudential concern, because it is about the fate of the system rather than merely the fate of one contracting party.

The lecture also suggests a historical shift. Earlier regulation, in the mix that actually prevailed, leaned more heavily toward the microprudential side. The crisis of the 2000s forced a stronger macroprudential emphasis.

\subsection{Question \& Answer}

\paragraph{Question.} Why can rescuing a large firm create more risk rather than less?

\paragraph{Answer.} Because the rescue changes incentives before the failure occurs. Once market participants understand that sufficiently large firms will be rescued to prevent broader collapse, those firms acquire a kind of free insurance. The taxpayer is not paid an actuarially fair premium for underwriting this insurance, but the firm can behave as if the insurance were there. That is why the lecture treats \emph{too big to fail} as a regulatory problem in its own right: the rescue is stabilizing in the short run, but destabilizing if it becomes built into expectations.

\section{Referees, Competition, and the Need for Rules}

Before he begins the institutional survey, the lecturer pauses to justify regulation itself. This pause is important. He does not proceed as if the value of rules were obvious. Instead he compares regulators to referees at a sports event.

The comparison is direct. Referees enforce rules; they decide when the rules have been broken; they can punish violations. Players may complain about referees, but they still want them, because without them dangerous play, strategic fouling, and deliberate injury would become close to necessary for success.

That same logic is transferred to finance. In a competitive system, if everyone else is doing something shady, we may feel pressured to do it too. The lecture calls this a kind of race to the bottom. Once that logic is in place, the demand for regulation no longer looks like an external political imposition; it can emerge from market participants themselves, who want somebody to keep the game from degenerating.

This defense of regulation also carries a moral point. We celebrate players more than referees, and in finance we tend to celebrate dealmakers more than regulators. But the lecture insists that regulation is a profession, and sometimes a morally attractive one. Some people move from Wall Street to regulatory work not because they failed, but because the rewards of life are plural and because they would rather enforce decent rules than operate permanently in a compromised environment.

\subsection{Question \& Answer}

\paragraph{Question.} If competition is supposed to discipline behavior, why do competitors still ask for regulation?

\paragraph{Answer.} Because competition disciplines some margins and corrupts others. It may discipline price or performance, but if rule-breaking is profitable and widespread, then competition can force honest participants to imitate dishonest ones. A referee prevents the equilibrium from drifting toward conduct that everyone privately dislikes but no one can afford to renounce unilaterally.

\section{A Hierarchy of Regulatory Scope}

The lecture now reorganizes itself with a blackboard outline. On the left are the terms already introduced---\emph{Too big to fail}, \emph{Microprudential}, and \emph{Macroprudential}. On the right is a second organizing device: a five-level ladder of regulatory scope.

\begin{figure}[t]
\centering
\includegraphics[width=0.88\textwidth]{lecture_12_figure_02.png}
\caption{Blackboard outline of prudential regulation and the five levels of regulatory scope.}
\end{figure}

The important point is not that the blackboard gives a formal diagram. It does not. There are no arrows, braces, or connecting lines. What it does give us is the lecture's architecture. The prudential distinction has already been introduced, and now the lecturer places beside it a hierarchy of institutional and geographical scope:
\begin{enumerate}
\item Within firm
\item Trade groups
\item Local government
\item National
\item International
\end{enumerate}

A clean typographic reconstruction of the board vocabulary is
\begin{align}
\text{Too big to fail},\qquad
\text{Microprudential},\qquad
\text{Macroprudential}.
\end{align}

The lecture immediately comments on the historical tendency behind this ladder. The older forms of regulation are the narrower ones. Over time regulation has become broader in geographical scope, which is why international regulation is saved for the end of the lecture.

\section{Within-Firm Regulation: Boards, Tunneling, and Fiduciary Duty}

The first level is within the firm. The lecturer explicitly counts this as regulation even though it does not begin with the government. Firms impose standards on themselves, and in the modern corporation the principal internal body is the board of directors. He even says that the board could be called a board of regulators, though no one uses that phrase.

The quasi-regulatory character of the board comes mainly from the outside directors. Inside directors work for the firm; outside directors bring independence, reputation, and community standing. The lecture leans heavily on the social meaning of this. A board populated by widely respected figures sends a signal that certain kinds of conduct will not be tolerated.

This is why the Yale Corporation appears as an example. The names are not mathematically important, but institutionally they are. The governor of the state, well-known public intellectuals, major corporate leaders, and figures who serve on other prominent boards all contribute reputational capital. Interlocking directorships matter because the quality of a director is judged in part by the other boards on which that person is willing to serve.

But the lecture then sharpens the argument. Why does a for-profit corporation need this internal regulation? Because there are too many opportunities for insiders to divert value away from the shareholders. This is the problem the lecturer calls \emph{tunneling}. Tunneling means extracting value from the corporation for insiders rather than leaving that value with the shareholders who supplied the capital.

The lecture's catalogue of tunneling channels is worth preserving:
\begin{itemize}
\item \textbf{Asset sales.} An insider arranges to sell a corporate asset at a non-market price to an associate, who later returns part of the gain privately.
\item \textbf{Contracts.} A firm overpays a favored supplier in exchange for later kickbacks or reciprocal favors.
\item \textbf{Executive compensation.} A chief executive pays friends and cronies inflated salaries and defends them as justified.
\item \textbf{Expropriation of corporate opportunities.} A valuable business opportunity that properly belongs to the corporation is diverted into a new private vehicle controlled by insiders.
\item \textbf{Insider trading.} Managers trade on nonpublic information generated by their role inside the firm.
\end{itemize}

The logic is cumulative. The lecture does not mention these as isolated legal curiosities. It piles them up so that we feel the problem: if insiders control so many margins of action, why should outside investors trust the corporation at all?

\subsection{Question \& Answer}

\paragraph{Question.} Why would anyone invest in a corporation if insiders have so many ways to tunnel value out of it?

\paragraph{Answer.} Because the corporation is not only a productive organization; it is also a governance structure. The board, especially its outside members, exists to monitor management, to impose standards of conduct, and to put personal and public reputation behind the claim that the firm's affairs are being run honestly. Without that layer of internal regulation, the corporation would indeed be too delicate a form to attract outside capital.

The lecture then formulates the board member's obligations in traditional legal language. First is the \emph{duty of care}: the duty to act as a reasonable, prudent, rational person would act. In the lecture this means paying attention, getting information, and monitoring management rather than passively accepting presentations. A board member does not manage the company line by line, but must know enough to function as a regulator.

Second is the \emph{duty of loyalty}. The lecture notes that there is increasing discussion of wider obligations to the community, but its baseline remains classical: loyalty is principally owed to the shareholders, because they put up the money and are the parties most exposed to tunneling.

Thus the within-firm level already contains a complete regulatory logic:
\begin{enumerate}
\item shareholders supply capital;
\item insiders control information and assets;
\item control creates channels for tunneling;
\item boards and fiduciary duties are introduced to block those channels.
\end{enumerate}

\section{Trade Groups, Local Regulation, and National Disclosure}

The lecture then ``moves up'' from the firm to trade groups. These are groups of firms that organize themselves into associations or self-regulatory organizations. The central example is the New York Stock Exchange, founded in 1792 shortly after the first American stock market crash.

The founding document is the Buttonwood Agreement. The lecture reads it because the document has often been romanticized, whereas its economic content is strikingly compact. In effect, the agreement does two things: it sets a minimum commission and it gives preference to trading among the members themselves. The commission term can be summarized as
\begin{equation}
c_{\text{Buttonwood}}=\frac{1}{4}\%=0.25\%=0.0025.
\end{equation}

In that narrow sense, the agreement looks like cartelization. Yet the lecturer does not stop there. He gives the more charitable interpretation: after a speculative crash driven by a dubious promoter, the brokers were also trying to exclude bad actors and maintain a ``gentleman's organization'' with enough revenue to support responsible conduct. This is the lecture's tension. A trade group may impose ethical standards and still simultaneously create monopoly profits.

\subsection{Question \& Answer}

\paragraph{Question.} How can self-regulation help markets without simply becoming a cartel?

\paragraph{Answer.} The lecture's answer is historical rather than axiomatic. Self-regulation can help by screening participants, policing conduct, and building reputational standards. But the same institutional form can also fix prices and exclude competition. Public regulation then has a second-order task: preserve surveillance and ethical rulemaking where they are useful, while breaking cartel power where it becomes monopolistic.

That second-order task appears in the U.S.\ reform known as \emph{May Day}, on May 1, 1975. Fixed stock-exchange commissions became illegal, and the closed, member-only commission structure inherited from the Buttonwood world was broken. The lecture's key point is not that the exchange ceased to exist; it is that its monopoly structure was dismantled. The same process later appeared in the United Kingdom's 1986 \emph{Big Bang}. In both cases the system moved toward competitive commissions while retaining the idea that surveillance against manipulation still mattered.

From trade groups the lecture moves to local regulation. In the United States, before the 1930s, financial regulation was overwhelmingly local. State-level \emph{Blue Sky Laws}, beginning in Kansas in 1911 and spreading during the Progressive Era, were designed to block fraudulent or absurdly overpromoted securities. But the limits of local regulation became obvious once finance crossed state lines. Boiler rooms, telephone banks, and cross-state solicitation made it easy to evade fragmented enforcement.

This failure of local scale leads directly to national regulation. The key institution is the Securities and Exchange Commission, founded in 1934 under the New Deal. The lecture presents the SEC as a disclosure-centered regulator. Louis Brandeis's famous phrase about sunshine being the best disinfectant is invoked to capture the idea: financial dishonesty thrives in secrecy, so public disclosure is itself a regulatory instrument.

The national-level logic is therefore:
\begin{enumerate}
\item local rulemaking is fragmented;
\item cross-border solicitation defeats fragmented enforcement;
\item a national regulator is created;
\item disclosure and surveillance become central tools.
\end{enumerate}

The lecture's treatment of the public-private distinction in securities fits naturally into this disclosure framework. A public company entering the market by an IPO must pass through a demanding regulatory procedure. A private fund, by contrast, is permitted more freedom, but only because access is restricted. The lecture describes two common private-fund screens in approximate quantitative form:
\begin{align}
N_{\text{class 1}} &\le 99,
&
y_{\text{single}} &\ge 200{,}000,
&
y_{\text{married}} &\ge 300{,}000,
&
W_{\text{investable}} &\ge 1{,}000{,}000,
\\
N_{\text{class 2}} &\le 500,
&
W_{\text{individual}} &\ge 5{,}000{,}000,
&
W_{\text{institution}} &\ge 25{,}000{,}000.
\end{align}

The lecture then pauses over the famous hedge-fund fee structure ``2 and 20.'' In a compact notation consistent with the discussion, if \(A\) denotes assets under management and \(\Pi\) trading profits, then
\begin{equation}
f(A,\Pi)=0.02A+0.20\max\{\Pi,0\}.
\end{equation}

\paragraph{Worked example.} If the fund manages one billion dollars, then even before any performance fee we have
\begin{equation}
f(10^9,\Pi)\ge 0.02\times 10^9 = 2\times 10^7 = 20{,}000{,}000.
\end{equation}
That is the lecturer's arithmetic point: the management fee alone is already large, and the performance fee sits on top of it while losses do not symmetrically fall on the manager.

This worked example is not offered in the lecture to build a theory of optimal contracting. It is offered to dramatize the two-tier structure of investor protection. Public investors are kept away from a domain that the lecturer describes as something like a jungle. Wealth screens are used as a substitute for ordinary regulatory protection.

The national section ends by broadening the notion of surveillance. The SEC does not work only through disclosure. It also works through market surveillance, often in conjunction with self-regulatory organizations such as the exchanges. The lecture's insider-trading and manipulation stories---the Lotus case and the Emulex case---both serve the same pedagogical purpose: without referees, the market can be captured by people who use privileged information or false information to take money from others.

Two additional national-level institutions deserve brief preservation because they fit the lecture's logic of delegated and protective regulation. FASB is a private body recognized as the arbiter of accounting standards, and GAAP is the framework within which public-company accounting is presented. SIPC, created in 1970 after brokerage failure exposed the vulnerability of investors holding securities in street name, plays for brokerage accounts a role analogous to what deposit insurance plays for bank accounts.

\section{Post-Crisis Reform and the International Turn}

The final movement of the lecture scales regulation outward again. In the United States the important post-crisis legislation is Dodd--Frank, enacted in 2010. The lecturer reads it as a shift toward the macroprudential side. Earlier laws were already aimed at fraud against individuals, though not always successfully. Dodd--Frank adds institutions whose purpose is explicitly systemic.

The most important example is the Financial Stability Oversight Council, FSOC. Its task is to worry about systemic risk and the \emph{too big to fail} problem. Alongside it, the lecture places the Bureau of Consumer Financial Protection as a microprudential counterweight: one body focused on the system as a whole, another focused on ordinary borrowers and consumers.

The same two-level structure then appears in Europe. The European Supervisory Framework, also associated with 2010 reforms, includes a systemic body and sectoral bodies:
\begin{itemize}
\item the European Systemic Risk Board in Frankfurt;
\item the European Banking Authority in London;
\item the European Securities Market Authority in Paris;
\item the European Insurance and Occupational Pension Authority in Frankfurt.
\end{itemize}

The lecture presents these as signs that regulation is being lifted from the level of the individual country to a wider economic space.

That widening is then pushed to the international level. The problem is straightforward. If regulation is costly, financial firms can move offshore to jurisdictions with lighter requirements. Once that is possible, purely national regulation becomes unstable. The lecture therefore turns to international coordination bodies.

A brief chronology helps:
\begin{align}
1792 &: \text{Buttonwood Agreement and the NYSE},\\
1911 &: \text{Kansas Blue Sky law},\\
1934 &: \text{SEC},\\
1970 &: \text{SIPC},\\
1975 &: \text{May Day},\\
1986 &: \text{Big Bang},\\
2008 &: \text{G20 expansion},\\
2010 &: \text{Dodd--Frank, European supervisory reforms, Basel III approval}.
\end{align}

The Bank for International Settlements, in Basel, is presented as an old international coordinating institution whose influence comes not from direct legal command but from repeated consultation. The BIS is described as involving roughly 57 central banks. Closely related is the Basel Committee, whose major phases are
\begin{equation}
\text{Basel I }(1988),\qquad \text{Basel II }(2004),\qquad \text{Basel III }(2010).
\end{equation}

The lecture's tone here is careful. These are not global sovereigns. Their authority weakens as scope broadens. International regulation often works by recommendation, coordination, and mutual pressure rather than by direct coercion. But the lecture does not treat that weakness as insignificance. Quite the contrary: after a worldwide crisis, even nonbinding coordination can matter greatly.

The same movement is visible in the path from the G6 to the G7 and then to the G20. The G20, in particular, is described as a broader and more representative forum, and the Financial Stability Board is presented as one of the main channels through which global regulatory recommendations are now formulated and transmitted.

At the end, the lesson is historical and structural rather than merely institutional. Finance has become more interconnected and more global. For that reason regulation has had to become more macroprudential and more international.

\section{Summary}

The lecture begins with human misbehavior and ends with world-scale institutions. In between, it develops a consistent line of thought. Regulation is needed because finance is not only a mechanism for allocating capital; it is also a setting in which incentives, secrecy, reputation, monopoly, and fraud matter. At the system level, \emph{too big to fail} creates moral hazard and therefore demands macroprudential thinking. At the firm level, tunneling threatens the viability of the corporation and therefore demands boards, duties, and internal standards. At the industry level, self-regulation can improve conduct but also become cartelization. At the national level, disclosure and surveillance respond to the failure of local fragmentation. And at the international level, coordination becomes necessary because capital can move where regulation is lightest.

The lecture's closing claim is therefore not incidental. Financial regulation is becoming broader in scale because the financial system itself has become broader in scale. After global crisis, the problem can no longer be treated as local misconduct alone; it has to be understood as a problem of system design.